
\subsubsection{6.1 Interpretation of the PARTIAL Gate Result}

The experiment satisfies two of three gate criteria (ratio and AUC) by large margins. The ratio of 8.805 is nearly three times the 3.0x target; AUC = 0.914 is 0.214 above the 0.70 target. The third criterion (balance\_deviation $\leq$ 0.10) is not satisfied, but post-hoc analysis demonstrates that the $\leq$ 0.10 target is incompatible with the dataset's class imbalance and the definition of the criterion. The reflection decision recorded in \texttt{h-e1/reflection\textbackslash \_report.md} is SELF\_MODIFY: the criterion is to be replaced with minority recall $\geq$ 0.60 in the next iteration (h-e1-v2), which was not yet executed.

The AUC = 0.914 and ratio = 8.805 results are obtained from a single experimental run (seed = 42). Multi-seed replication (5 seeds) was planned as hypothesis h-m4 but was not executed. These values should therefore be interpreted as single-seed results pending confirmation.

\subsubsection{6.2 Mechanistic Account}

The outer-product decomposition $\|\nabla _W \ell _{i}\|_F = \|p_i - y_{i}\| \cdot \|h(x_{i})\|$ makes a clear mechanistic prediction: if BatchNorm equalizes $\|h(x_{i})\|$ across groups (confirmed, h\_norm\_std\_ratio $\approx$ 0.10), then the gradient norm disparity between groups must reflect prediction residual differences. The observed 8.8x ratio between minority and majority $\tilde{g}$ values at epoch 5 is consistent with this prediction.

However, the mechanistic account makes no claim about the magnitude of the effect, only its direction. The empirically observed 8.8x ratio could be influenced by factors beyond the outer-product decomposition alone --- including specific properties of the Waterbirds dataset, the ImageNet pretraining of the ResNet-50 backbone, and the particular training configuration used. Generalization to other datasets (e.g., CelebA) and architectures was not tested.

\subsubsection{6.3 Relationship to EL2N}

The normalized gradient norm $\tilde{g}_i = \|p_i - y_{i}\|$ is algebraically equivalent to the EL2N score \citep{Paul2021} when evaluated at a single training checkpoint rather than averaged over early training. EL2N was proposed for dataset pruning --- identifying and removing easy (low-error) samples --- while the present work uses the same quantity to identify hard (high-error) minority samples for a different downstream task. The two use cases are complementary in the selection direction (bottom-k\% vs. top-k\%) and distinct in purpose (pruning vs. minority enrichment for retraining).

\subsubsection{6.4 Limitations}

\textbf{Limitation 1: WGA not measured.} The primary limitation of this work is that Stage 2 (last-layer retraining on the $\tilde{g}-selected$ subset) was not executed. There are no worst-group accuracy results for the proposed GNR-LLR pipeline. Whether the high AUC of the proxy signal translates to improved WGA is an open empirical question.

\textbf{Limitation 2: Single dataset.} All experiments were conducted on Waterbirds only. No CelebA experiments were run. Generalization of the signal to other spuriously correlated benchmarks is unknown.

\textbf{Limitation 3: Minority recall not directly measured.} AUC = 0.914 provides indirect evidence that top-k\% selection captures a large fraction of minority samples, but the direct minority recall at k = 25\% was not computed in H-E1. A revised experiment (h-e1-v2) with minority recall as the gate criterion was proposed but not yet run.

\textbf{Limitation 4: Single random seed.} All results use seed = 42. Multi-seed validation was not performed. Results may vary across seeds.

\textbf{Limitation 5: Temporal behavior beyond epoch 10 not measured.} The ratio plateau at epoch 10 (8.5x) is consistent with majority saturation but this interpretation was not tested by collecting $\tilde{g}$ beyond epoch 10.

